% Part: first-order-logic
% Chapter: model-theory
% Section: nonstandard-arithmetic

\documentclass[../../../include/open-logic-section]{subfiles}

\begin{document}

\olfileid{mod}{bas}{nsa}
\section{算术的非标准模型}

\begin{defn}
设 $\Lang{L_N}$ 为算术语言，包含常项符号 $\Obj{0}$、二元谓词符号 $<$、一元函数符号 $\prime$ 以及二元函数符号 $+$ 和 $\times$。
\begin{enumerate}
\item 算术的\emph{标准模型}是 $\Lang{L_N}$ 的结构 $\Struct{N}$，其论域为 $\Nat = \{ 0, 1, 2, \dots\}$
且将 $\Obj{0}$ 解释为 $0$，将 $<$ 解释为 $\Nat$ 上的小于关系，将 $\prime$、$+$ 和 $\times$ 分别解释为 $\Nat$ 上的后继、加法和乘法。
\item \emph{真算术}即理论 $\Theory{N}$。
\end{enumerate}
\end{defn}

在 $\Lang{L_N}$ 中工作时，我们将形如 $\Obj{0}^{\prime\dots\prime}$（即对 $\Obj{0}$ 应用 $n$ 次后继函数）的项简记为~$\num{n}$。

\begin{defn}
一个 $\Lang{L_N}$ 的结构~$\Struct{M}$ 是\emph{标准的}，当且仅当 $\Struct{N} \iso \Struct{M}$。
\end{defn}

\begin{thm}\ollabel{thm:non-std}
真算术存在非标准的可数模型。
\end{thm}

\begin{proof}
扩充语言 $\Lang{L_N}$，引入一个新常项符号~$c$，并考虑理论
\[
\Theory{N} \cup \Setabs{\num{n} < c}{n \in \Nat}.
\]
该理论是有限可满足的，因此由紧致性定理可知它有一个模型 $\Struct{M}$。根据 Löwenheim–Skolem 下降定理，可以取 $\Struct{M}$ 为可数的。设 $\Domain{M}$ 为 $\Struct{M}$ 的论域，并设 $\Struct{M}$ 将 $\Lang{L}$ 的非逻辑常项解释为 $\mathbf{z} = \Assign{\Obj{0}}{M} \in \Domain{M}$、${\prec} =
\Assign{<}{M} \subseteq M^2$、$* = \Assign{\prime}{M} \colon
\Domain{M} \to \Domain{M}$ 和 $\oplus = \Assign{+}{M}, \otimes =
\Assign{\times}{M} \colon \Domain{M}^2 \to \Domain{M}$。对每个 $x \in \Domain{M}$，我们用 $x^*$ 表示对 $x$ 应用 $*$ 所得到的 $\Domain{M}$ 中的元素。

现在，如果 $h$ 是 $\Struct{N}$ 与 $\Struct{M}$ 之间的同构，则存在 $n \in \Nat$ 使得 $h(n) = \Assign{c}{M}$。于是，设 $s$ 为 $\Struct{N}$ 中满足 $s(x) = n$ 的任意赋值。则 $\Sat{N}{\eq[\num{n}][x]}[s]$；根据 \olref[iso]{thm:isom} 的证明，亦有 $\Sat{M}{\eq[\num{n}][x]}[h \circ s]$，从而 $\Assign{c}{M} =
\mathbf{z}^{*\cdots *}$（其中 $*$ 迭代了 $n$ 次）。但这是不可能的，因为由假设 $\Sat{M}{\num{n} < c}$ 且 $\prec$ 是反自反的。所以 $\Struct{M}$ 是非标准的。
\end{proof}

\begin{prob}
集合 $X$ 上的一个关系 $R$ 是\emph{良基的}，当且仅当 $R$ 中不存在无限下降链，即不存在 $X$ 中的元素 $x_0$, $x_1$, $x_2$, \dots 使得 $\dots x_2Rx_1Rx_0$。假设 Zermelo–Fraenkel 集合论 $ZF$ 是一致的，证明存在 $ZF$ 的\emph{非良基模型}，即满足 $\dots x_2 \in x_1 \in x_0$ 的模型 $\mathfrak{M}$。
\end{prob}

由于 \olref{thm:non-std} 中的非标准模型 $\Struct{M}$ 与标准模型初等等价，我们可以推导出 $\Struct{M}$ 的许多性质。本节余下部分将致力于此，以便精确刻画 $\Theory{N}$ 的可数非标准模型。

\begin{enumerate}
\item $\Domain{M}$中没有成员$\prec$小于自身：语句
  $\lforall[x][\lnot x < x]$在$\Struct{N}$中为真，从而在
  $\Struct{M}$中也为真。
\item 通过类似的推理，我们得到$\prec$是$\Domain{M}$的一个\emph{线性序}，
  即$\Domain{M}$上的一个全、反自反且传递的关系。
\item 元素$\mathbf{z}$是$\Domain{M}$的$\prec$最小元。
\item $\Domain{M}$的任何成员都$\prec$小于其$*$-后继，并且
  $x^*$是$\Domain{M}$中大于$x$的$\prec$最小成员。
\item $\Struct{M}$包含一个（关于$\prec$的）初始段，它同构于
  $\Nat$：$\mathbf{z}, \mathbf{z}^*, \mathbf{z}^{**}, \dots$，我们称之为
  $\Domain{M}$的\emph{标准部分}。$\Domain{M}$的任何其他成员
  都是\emph{非标准的}。$\Domain{M}$中必须有非标准成员，否则
  \olref{thm:non-std}证明中的函数$h$就是一个同构。
  我们用$n, m, \dots$作为变元，在此标准部分上取值。
\item 每个非标准元素都大于任何标准元素；
  这是因为对于每个$n \in \Nat$，
  \[
  \Sat{N}{\lforall[z][(\lnot(\eq[z][\Obj{0}] \lor
  \dots \lor \eq[z][\num{n}]) \lif \num{n} < z)]},
  \]
  所以，如果$z \in \Domain{M}$不同于所有标准元素，它必须\emph{大于}
  所有标准元素。
\item $\Domain{M}$中除$\mathbf{z}$外的任何成员，都是
  $\Domain{M}$中某个唯一元素的$*$-后继，记作$^*x$。
  如果$x = y^*$，那么若其中一个是标准的，则两者皆是标准的
  （若其中一个是非标准的，则两者皆是非标准的）。
\item 在$\Domain{M}$上定义等价关系$\approx$为：
  $x \approx y$当且仅当存在某个\emph{标准}的$n$，使得
  要么$x \oplus n = y$，要么$y \oplus n =x$。
  换言之，$x \approx y$当且仅当$x$和$y$相距有限距离。
  如果$n$和$m$都是标准的，那么$n \approx m$。
  定义$x$的\emph{块}为其等价类$[x] = \Setabs{y}{x
  \approx y}$。
\item 假设$x \prec y$且$x \not\approx y$。由于
  $\Sat{N}{\lforall[x][\lforall[y][(x < y \lif (x' < y \lor x' =
      y))]]}$，要么$x^* \prec y$，要么$x^* = y$。
  后者不可能，因为这会推出$x \approx y$，所以$x^* \prec
  y$。类似地，如果$x \prec y$且$x \not\approx y$，那么$x \prec
  {^*y}$。因此，如果$x \prec y$且$x \not\approx y$，那么每个
  与$x$等价的$w$都$\prec$小于每个与$y$等价的$v$。相应地，
  每个块$[x]$构成一个双向无穷链
  \[
  \cdots \prec  {^{**}x} \prec {^*}x \prec x \prec x^* \prec x^{**}
  \prec \cdots
  \]
  它被称为一个$Z$-链，因为其序型与整数集相同。
\item 可以将$\prec$序提升到块上：如果$x \prec y$，
  则$x$所在的块小于$y$所在的块。一个块如果包含非标准元素，它就是
  \emph{非标准的}。标准块是最小的块。
\item 不存在最小的非标准块：如果$y$是非标准的，
  那么存在$x \prec y$，其中$x$也是非标准的且$x
  \not\approx y$。证明：在标准模型$\Struct{N}$中，每个
  数都能被2整除（可能余1），即
  $\Sat{N}{\lforall[y][\lforall[x][(y = x + x \lor y = x + x +
        \Obj{0}')]]}$。根据初等等价，对于每个$y \in
  \Domain{M}$，都存在$x \in \Domain{M}$使得要么$x \oplus x
  = y$，要么$x \oplus x \oplus \mathbf{z}^*= y$。如果$x$是标准的，
  那么$y$也会是标准的；所以$x$是非标准的。此外，$x$和
  $y$属于不同的块，即$x \not\approx y$。要明白这一点，
  假设它们属于同一块，即对某个标准$n$有$x \oplus n =
  y$。如果$y = x \oplus x$，那么$x \oplus n = x
  \oplus x$，由此根据加法消去律（它在$\Struct{N}$中成立，因此在
  $\Struct{M}$中也成立），可得$x = n$，从而$x$仍是标准的。
  类似地，如果$y = x \oplus x
  \oplus \mathbf{z}^*$。
\item 根据类似的论证，不存在最大的块。
\item 块的序是稠密的：如果$[x]$小于$[y]$
  （其中$x \not\approx y$），那么存在一个不同于两者的块$[z]$介于它们之间。
  假设$x \prec y$。如前所述，$x \oplus
  y$能被2整除（可能余1），所以存在$u
  \in \Domain{M}$使得要么$x \oplus y = u \oplus u$，要么$x
  \oplus y = u \oplus u \oplus \mathbf{z}^*$。元素$u$是
  $x$与$y$的平均值，因此介于它们之间。假设$x \oplus y =
  u \oplus u$（另一种情况类似）：如果$u \approx x$，那么对某个标准$n$：
  \[
  x \oplus y = x \oplus n \oplus x \oplus n,
  \]
  于是$y = x \oplus n \oplus n$，则可得$x \approx y$，
  与假设矛盾。我们断定$u \not\approx x$。类似的
  论证可得出$u \not\approx y$。
\end{enumerate}

因此，非标准块的排序如同有理数集：它们构成一个没有端点的可数线性序。
由此可得，对于真算术的任意两个可数非标准模型 $\Struct{M}_1$ 和 $\Struct{M_2}$，
它们在仅含 $<$ 和 $=$ 的语言上的归约是同构的。
事实上，可以如下定义一个同构 $h$：
$\Struct{M}_1$ 和 $\Struct{M}_2$ 的标准部分都同构于标准模型 $\Struct{N}$，从而彼此同构。
构成非标准部分的那些块本身如同有理数集般排序，因此根据 \olref[dlo]{thm:cantorQ}，它们也是同构的；
块之间的同构可以通过以下方式扩展到块\emph{内部}的同构：先在每对对应的块中各任选一个元素进行匹配，然后将 $\Struct{M}_1$ 中 $x$ 的后继的像取为 $\Struct{M}_2$ 中 $x$ 的像的后继。
注意，这\emph{并}不\emph{意味着} $\mathfrak{M}_1$ 和 $\mathfrak{M}_2$ 在算术的完整语言中也是同构的（实际上，同构总是相对于一个符号集的），因为在 $\Domain{M_1}$ 和 $\Domain{M_2}$ 上定义加法和乘法存在非同构的方式。（这一点也可以从 Vaught 的一个著名定理得出：一个完备理论的可数模型个数不可能为~2。）

\begin{prob}
证明在算术的非标准模型中不可能存在最大的块。
\end{prob}

\begin{prob}
设 $\Lang{L}$ 是仅以 $<$ 为谓词符号（等词 $\eq$ 除外）的一阶语言，
并令 $\Struct{N} = (\Nat, <)$。$\Struct{N}$ 的所有有限或余有限子集都是可定义的。
证明这些是 $\Struct{N}$ 的\emph{仅有的}可定义子集。

（提示：首先，令 $\Obj{prc}(x,y)$ 为缩写“$x$ 是 $y$ 的直接前驱”的 $\Lang{L}$-公式：
\[
x<y \land \lnot \lexists[z][(x<z \land z < y)].
\]
现在，对于 $\Struct{N}$ 的任意可定义子集，都对应着 $\Lang{L}$ 中的一个公式 $!A(x)$。
对任意这样的 $!A$，考虑语句 $!D$：
\[
\lexists[x][\lforall[y][\lforall[z][((x<y \land x<z \land \Obj{prc}(y,z)
  \land !A(y)) \lif !A(z))]]].
\]
证明 $\Sat{N}{!D}$ 当且仅当由 $!A$ 定义的 $\Struct{N}$ 的子集是有限或余有限的。

现在，设 $\Struct{M}$ 是一个与 $\Struct{N}$ 初等等价的非标准模型。
如果 $a \in \Domain{M}$ 是非标准的，取 $b, c \in \Domain{M}$ 均大于 $a$，并且令 $b$ 是 $c$ 的直接前驱。
那么存在 $\Domain{M}$ 的一个自同构 $h$ 使得 $h(b)=c$（为什么？）。
因此，如果 $b$ 满足 $!A$，那么 $c$ 亦满足 $!A$（为什么？）。
由此可得 $!D$ 在 $\Struct{M}$ 中为真，从而在 $\Struct{N}$ 中亦真。
但这意味着由 $!A$ 定义的 $\Struct{N}$ 的子集是有限或余有限的。）
\end{prob}

\end{document}